\documentclass[interno]{gpt}

\titolo{Valutazione della documentazione per RTB}
\data{2026-11-05}
\orario{18:00}{20:00}
\luogo{Discord}
\redattore{cescobolla}
\verificatore{corro}
%\assenti{spagnoli}

\modifica{1.0.0}{2026-11-06}{cescobolla}{-}{Approvazione verbale}
\modifica{0.1.0}{2026-11-06}{cescobolla}{corro}{Prima stesura}

\componenti{biasini,cescobolla,corro,guglielmo,mazzurana,spagnoli}
\begin{document}

\section{Ordine del giorno}
\begin{enumerate}
  \item Prenotazione dell'incontro con il proponente Eggon tramite calendario condiviso.
  \item Valutazione di gruppo della documentazione necessaria per la milestone RTB.
  \item Collaudo e verifica del sistema di tracciamento orario integrato con le issue di GitHub.
  \item Ultimazione delle modifiche al sito web di progetto a seguito dell'assegnazione dell'appalto.
  \item Pianificazione del materiale per l'incontro con il proponente Eggon.
\end{enumerate}

\clearpage

\section{Svolgimento}

\subsection{Prenotazione dell'incontro con Eggon S.r.l.}
Il gruppo ha formalizzato la prenotazione del primo incontro con la proponente Eggon S.r.l., avvalendosi del calendario delle disponibilità condiviso dall'azienda via email. La riunione è stata fissata per mercoledì 11 novembre 2026 alle ore 16:15 su Google Meet.

\subsection{Valutazione congiunta della documentazione per la milestone RTB}
A seguito dell'attività di ricerca condotta, il team si è confrontato per condividere una comprensione omogenea della documentazione richiesta in vista della milestone RTB. La discussione ha permesso di allineare tutti i componenti sulla struttura dei nuovi documenti e di definire la linea operativa per la successiva stesura.

\subsection{Collaudo del tracciamento orario automatizzato tramite GitHub}
È stato effettuato con esito positivo il test per il conteggio orario automatizzato basato sulle issue di GitHub. Il sistema garantisce un tracciamento preciso e trasparente delle ore lavorative rendicontate da ciascun membro suddivise per ruolo.

\subsection{Ultimazione delle modifiche al sito web di progetto}
Il team ha verificato e completato le ultime modifiche al sito web di progetto a seguito delle criticità emersi in sede di aggiudicazione dell'appalto, aggiornando la struttura di visualizzazione e consultazione della documentazione esposta.

\subsection{Analisi dei requisiti principali e ripartizione dei task}
In vista del primo incontro con Eggon S.r.l., il gruppo ha stabilito di elaborare una prima analisi del capitolato d'appalto isolandone i requisiti principali. Si è quindi proceduto a ripartire le sezioni del capitolato tra i membri del team per avviare lo studio individuale.

\section{Decisioni}
\begin{decisioni}
    \decisione{Prenotazione e conferma dell'incontro con la proponente Eggon S.r.l. fissato a calendario per mercoledì 11 novembre 2026 alle ore 16:15.}
    \decisione{Adozione definitiva del sistema di conteggio orario automatizzato integrato nelle issue di GitHub.}
    \decisione{Completamento ed approvazione delle modifiche alla visualizzazione della documentazione sul sito web di progetto.}
    \decisione{Suddivisione delle sezioni del capitolato d'appalto tra i componenti per l'analisi individuale dei requisiti principali.}
\end{decisioni}

\section{Azioni da intraprendere}
\begin{azioni}
  \azione[\#84]{Stesura verbale interno della riunione odierna.}{cescobolla}{05/11/2026}
  \azione[\#87]{Analisi dei requisiti principali dell'appalto.}{tutti}{11/11/2026}
\end{azioni}

\section{Prossima riunione}
La prossima riunione interna è fissata per martedi 10 novembre 2026 alle ore 18.00 con lo scopo di revisionare e ultimare i contenuti da presentare al proponente Eggon durante il successivo incontro. 

\end{document}